
L'idéal engendré par \(P\) est un \(\lambda\)-idéal. Pour le voir, on remarque d'abord que les relations \(\lambda^i(N)=0\) pour \(i>d\) entraînent que les \(\lambda^j(N-L)\) pour \(j\ge d\) sont multiples de \(\lambda^d(N-L)\). Si on se place dans \(1+\widehat{K[L]/(P)[[t]]}^+\), l'image de \(\lambda_t(N-L)\) est un polynôme de degré \(<d\), et un argument universel montre alors que ses \(\lambda^i\) pour \(i\ge d\) sont égaux à \(1\) (cf. V 4.8). Comme dans \(1+\widehat{K[L][[t]]}^+\) on a \(\lambda_t(\lambda^d(N-L))=\lambda^d(\lambda_t(N-L))\), on en déduit que les \(\lambda^i(\lambda^d(N-L))\) pour \(i\ge1\) sont multiples de \(P\). D'après les relations V 2.4.2, cela entraîne que l'idéal engendré par \(P\) est un \(\lambda\)-idéal.

On pose alors \(K_1=K[L]/(P)\). L'homomorphisme naturel \(K\to K_1\) sera noté \(g^*\); c'est un \(\lambda\)-homomorphisme. D'autre part, \(K_1\) est un \(K\)-module libre ayant pour base les éléments \(1,L,\ldots,L^{d-1}\). Il résulte de la relation \(\lambda^d(N)=\lambda^{d-1}(N-L).L\) et de l'hypothèse d'inversibilité faite sur \(\lambda^d(N)\) que \(L\) est inversible. Par suite \(K_1\) admet pour base sur \(K\) les éléments \(1,L^{-1},\ldots,L^{-d+1}\) ; on définit alors une application \(g_*:K_1\to K\) en associant à tout élément \(y'=y_0+y_1L^{-1}+\ldots+y_{d-1}L^{-d+1}\) de \(K_1\) l'élément \(y_0\) de \(K\). Il est clair que l'application ainsi définie est un homomorphisme additif, \(K\)-linéaire, et tel que \(g_*\circ g^*=\Id\).

On munit ensuite \(K_1\times K'\) d'une structure de \(\lambda\)-anneau, en posant :
\begin{align}
 (y'_1,x_1)*(y'_2,x_2)
   &=(y'_1y'_2(1-L)+x_1y'_2+x_2y'_1,\; x_1x_2); \tag{4.7.1}\\
 \lambda^n_*(y',x)
   &=\left(\sum_{p=1}^{n}\lambda^{\,n-p}(x)\lambda^p(L,y'),\;
     \lambda^n(x)\right). \tag{4.7.2}
\end{align}
Dans ces formules, \(K_1\) est considéré comme \(K'\)-algèbre par l'homomorphisme \(g^*\circ i^*\).
On pose ensuite, pour \(y_1,y_2\in K\) et \(n\ge1\) :
\begin{align}
 y_1*y_2&=-\,y_1y_2\lambda_{-1}(N); \tag{4.7.3}\\
 \lambda^n_*(y_1)&=-\,\lambda^n(N,-y_1); \tag{4.7.4}
\end{align}
